% Source-aligned English for printed p.170.
\SourcePage{170}

\LayerTarget{tohoku:proof:3.6.2:p2}
\noindent associated with an injective resolution of $A_i$, we have
$H^pT(C_i)=R^pT(A_i)$, and hence the desired conclusion

\LayerTarget{tohoku:eq:derived-colimit-isomorphism}
\[
R^pT(A)=\varinjlim_i R^pT(A_i).
\]

\LayerTarget{tohoku:proposition:3.6.3}
\noindent\textsc{Proposition 3.6.3.} --- \textit{Let $X$ be a topological
space endowed with an antifilter $\Phi$ of closed subsets. Then the functors
$H^p_\Phi(X,F)$ on $C^X$ commute with inductive limits in the following two
cases:}
\begin{enumerate}
\LayerTarget{tohoku:subitem:3.6.3:1}
\item \textit{$X$ is locally compact and $\Phi$ is the family of compact
subsets.}
\LayerTarget{tohoku:subitem:3.6.3:2}
\item \textit{$X$ is a Zariski space and $\Phi$ is the family of all its
closed subsets.}
\end{enumerate}

\LayerTarget{tohoku:def:zariski-space}
(We call a space a \textit{Zariski space} if every decreasing sequence of
closed subsets is stationary; cf. [15, page 223].) ---

\LayerTarget{tohoku:proof:3.6.3}
It suffices to verify conditions a) and b) of Proposition 3.6.2 for the functor
$\Gamma_\Phi$. Verifying a) is an easy compactness exercise left to the reader
(see also [9]). Condition b) follows from the corollary to Proposition 3.3.1
and from the following lemma.

\LayerTarget{tohoku:lemma:3.6.4}
\noindent\textsc{Lemma 3.6.4.} --- \textit{Under the assumptions of 1), every
inductive limit of $\Phi$-soft sheaves is $\Phi$-soft. Under the assumptions of
2), every inductive limit of flabby sheaves is flabby.}

\LayerTarget{tohoku:proof:3.6.4}
Assume the conditions of 1). Let $(F_i)$ be an inductive system of
$\Phi$-soft sheaves on $X$, let $F$ be its inductive limit, and let $f$ be a
section of $F$ over some $A\in\Phi$, that is, over a compact subset $A$. By a)
applied to the compact space $A$, the section $f$ comes from a section $f_i$ of
some $F_i$ over $A$. Since $F_i$ is $\Phi$-soft, $f_i$ is the restriction of
some $g_i\in\Gamma_\Phi(F_i)$; thus $f$ is the restriction of the section
$g\in\Gamma_\Phi(F)$ defined by $g_i$. Case 2) is entirely analogous, since a
subset $U$ of a Zariski space is itself a Zariski space and a) therefore
applies to it.

\LayerTarget{tohoku:def:cohomological-dimension}
We say that a space $X$ has \textit{cohomological dimension} $\leq n$ if
$H^i(X,F)=0$ for $i>n$ and every sheaf of abelian groups $F$ on $X$.
Combining Propositions 3.6.1 and 3.6.3 gives the following result.

\LayerTarget{tohoku:corollary:cohomological-dimension-test}
\noindent\textsc{Corollary.} --- \textit{Let $X$ be either a compact space or
a Zariski space, and let $n$ be an integer $\geq0$. Then $X$ has
cohomological dimension $\leq n$ if and only if $H^i(X,Z_U)=0$ for $i>n$ and
every open subset $U$ of $X$.}

\LayerTarget{tohoku:def:combinatorial-dimension}
We now come to the essential result of this section. A Zariski space $X$ is
said to have \textit{combinatorial dimension} $\leq n$ if every strictly
decreasing sequence of irreducible closed subsets has at most $n+1$ members.
Then:

\LayerTarget{tohoku:theorem:3.6.5}
\noindent\textsc{Theorem 3.6.5.} --- \textit{Let $X$ be a Zariski space of
combinatorial dimension $\leq n$. Then $X$ has cohomological dimension
$\leq n$; that is, $H^i(X,F)=0$ for $i>n$ and every sheaf of abelian groups
$F$ on $X$.}

\LayerTarget{tohoku:proof:3.6.5:p1}
We argue by induction on the combinatorial dimension $n$ of $X$. The theorem
is trivial for $n=0$, when $X$ is a finite discrete set. Suppose it proved for
combinatorial dimensions $\leq n-1$, where $n\geq1$, and prove it when $X$ has
combinatorial dimension $\leq n$. Let the $X_i$ be the irreducible

\TohokuContinuation{171}
